%%%%%%%%%%%%%%%%%%%%%%%%%%%%%%%%%%
\makeatletter

\usepackage{iftex}

\usepackage[
	colorinlistoftodos,%
	% disable,
	color=thillux-grey!35,%
	bordercolor=thillux-grey,%
	linecolor=thillux-grey%
]{todonotes}

\ifxetex
\usepackage{fontspec}
\defaultfontfeatures{Ligatures=Common}
%\setmainfont{Source Serif Pro}
%\setmonofont[Scale=MatchLowercase]{Source Code Pro}
%\setsansfont[Scale=MatchLowercase]{Source Sans Pro}
\usepackage{amsmath}
\usepackage{amssymb}
\usepackage{unicode-math}
%\setmathfont[Ligatures=TeX]{Latin Modern Math}
\else
\ifluatex
\usepackage{fontspec}
\defaultfontfeatures{Ligatures=Common}
%\setmainfont{Source Serif Pro}
%\setmonofont[Scale=MatchLowercase]{Source Code Pro}
%\setsansfont[Scale=MatchLowercase]{Source Sans Pro}
\usepackage{amsmath}
\usepackage{amssymb}
\usepackage{unicode-math}
%\setmathfont[Ligatures=TeX]{Latin Modern Math}
\else
\usepackage[utf8]{inputenc}
\usepackage[T1]{fontenc}
\usepackage{lmodern}
\usepackage{amsmath}
\usepackage{amssymb}
\fi
\fi

\usepackage{mathtools}
\usepackage{amsthm}

\usepackage{siunitx}

\usepackage[english]{babel}
\usepackage{enumitem}

\usepackage{bm}
\usepackage[nice]{nicefrac}     % schräge Bruchstriche
\usepackage{algorithm}
\usepackage{algpseudocode}

\floatname{algorithm}{Algorithm}
\renewcommand{\listalgorithmname}{List of algorithms}

%%%%%%%%%%%%%%%%%%%%%%%%%%%%%%%%%%
\usepackage[hidelinks]{hyperref}
\usepackage{bookmark}                           % Einstellungen für hyperref
\usepackage{booktabs}                           % nice tables
\usepackage[hang,flushmargin]{footmisc}         % schönere Fußzeile
\usepackage{graphicx}                           % Einbinden von Abbildungen
\usepackage{xcolor}                             % Farben
\usepackage{lipsum}
\usepackage{blindtext}
\usepackage{biblatex}
\usepackage{wrapfig}
\usepackage{texnames}
\usepackage{subcaption}
%\usepackage{ccicons} % Creative Commons Icons
%\usepackage{nameref} % ermöglicht Referenzierung von Abschnittstiteln
\usepackage[autostyle=true,english=american]{csquotes}
\usepackage{microtype}
\usepackage{tabularx}
\usepackage{relsize}
\usepackage{ifthen}

\usepackage{tikz}
\usepackage{pgfplots}
\pgfplotsset{compat=1.17}
\usepgfplotslibrary{
	fillbetween,
	statistics,
}

\addbibresource{doc/bib/books.bib}
\addbibresource{doc/bib/papers.bib}
\addbibresource{doc/bib/rfcs.bib}
\addbibresource{doc/bib/websources.bib}

\renewcommand{\bibname}{Literature list}

\allowdisplaybreaks

%%%%%%%%%%%%%%%%%%%%%%%%%%%%%%%%%%

\DeclareMathOperator*{\argmin}{arg\,min}
\DeclareMathOperator*{\argmax}{arg\,max}

\usepackage{esvect}

%%%%%%%%%%%%%%%%%%%%%%%%%%%%%%%%%%

\usepackage{multicol}
\setlength{\columnsep}{1em}

%%%%%%%%%%%%%%%%%%%%%%%%%%%%%%%%%%

% http://html5up.net/read-only
\definecolor{thillux-green}{HTML}{4ACAA8}

\definecolor{thillux-red}{HTML}{E7746F}

\definecolor{thillux-grey}{HTML}{989898}

\definecolor{thillux-blue}{HTML}{0E94EC}

\usepackage{tikz}
\usetikzlibrary{shapes,decorations}

\newcommand{\colorboxedsec}[2]{
	\tikz{\node[headings/#1]{#2};}}

\newcommand*{\ie}{i.\,e.,~}
\newcommand*{\eg}{e.\,g.,~}

\newcommand*{\authors}[1]{\textsc{#1} et\,al.}

%%%%%%%%%%%%%%%%%%%%%%%%%%%%%%%%%%
%\newcounter{anforderung}
%\newcommand{\anfl}[1]{%
%	\refstepcounter{anforderung}%
%	(A\arabic{anforderung})%
%	\label{anf:#1}%
%}

\newcounter{assumption}
\setcounter{assumption}{0}
\newcounter{nfrequirement}
\setcounter{nfrequirement}{0}
\newcounter{frequirement}
\setcounter{frequirement}{0}
\newcounter{secrequirement}
\setcounter{secrequirement}{0}

\newcommand*{\mmid}{\,\middle|\,}

\renewcommand*{\theassumption}{(A\arabic{assumption})}
\renewcommand*{\thenfrequirement}{(NF\arabic{nfrequirement})}
\renewcommand*{\thefrequirement}{(F\arabic{frequirement})}
\renewcommand*{\thesecrequirement}{(S\arabic{secrequirement})}

\newcommand*{\assumptionautorefname}{Assumption}
\newcommand*{\nfrequirementautorefname}{Requirement}
\newcommand*{\frequirementautorefname}{Requirement}
\newcommand*{\secrequirementautorefname}{Requirement}

\newcommand{\assumption}[1]{%
	\refstepcounter{assumption}%
	\item[\theassumption] {#1}%
}

\newcommand{\functionalrequirement}[1]{%
	\refstepcounter{frequirement}%
	\item[\thefrequirement] {#1}%
}
\newcommand{\nonfunctionalrequirement}[1]{%
	\refstepcounter{nfrequirement}%
	\item[\thenfrequirement] {#1}%
}
\newcommand{\secrequirement}[1]{%
	\refstepcounter{secrequirement}%
	\item[\thesecrequirement] {#1}%
}

\newenvironment{functionalrequirements}{%
	\begin{description}%
}{%
	\end{description}%
}
\newenvironment{nonfunctionalrequirements}{%
	\begin{description}%
}{%
	\end{description}%
}
\newenvironment{secrequirements}{%
	\begin{description}%
}{%
	\end{description}%
}

\usepackage[toc,page]{appendix}

%%%%%%%%%%%%%%%%%%%%%%%%%%%%%%%%%%%%%%%%%%%%%%%%

% example boxes

\usetikzlibrary{calc}

\newcommand{\TXExample}[1]{
	\begin{center}%
		\begin{tikzpicture}%
			\node [draw=none, fill=thillux-grey!15, very thick,
			rectangle, inner sep=10pt, inner ysep=8pt] (box){%
				\begin{minipage}{0.85\textwidth}%
					#1%
				\end{minipage}%
				};%
				\draw[line width=0.5ex, thillux-grey] ($(box.south west) + (0,0.6pt)$) |- ($(box.north west) - (0,0.6pt)$);%
		\end{tikzpicture}%
	\end{center}%
}

\newcommand{\TXExampleWithTitle}[2]{
	\begin{center}%
		\begin{tikzpicture}%
			\node [draw=none, fill=thillux-grey!15, very thick,
			rectangle, inner sep=10pt, inner ysep=8pt] (box){%
				\begin{minipage}{0.85\textwidth}%
					\textbf{#1}\\%
					#2%
				\end{minipage}%
				};%
				\draw[line width=0.5ex, thillux-grey] ($(box.south west) + (0,0.6pt)$) |- ($(box.north west) - (0,0.6pt)$);%
		\end{tikzpicture}%
	\end{center}%
}

%%%%%%%%%%%%%%%%%%%%%%%%%%%%%%%%%%%%%%%%%%%%%%%%%%

\graphicspath{{doc/img/}}
\usepackage[font=small, labelfont=bf]{caption}
% #1 - options for includegraphics (optional)
% #2 - path to figure
% #3 - caption
% #4 - label
\newcommand{\fig}[4][]{
	\begin{figure}[!htp]
		\centering
		\includegraphics[#1]{#2}
		\caption{#3}
		\label{#4}
	\end{figure}
}

\newcommand{\figMinipage}[3][]{
    \begin{minipage}{#3}
        \centering
        \includegraphics[#1]{#2}
    \end{minipage}
}

\newenvironment{figsMinipage}[2]{%
    \def\figsCaption{#1}%
    \def\figsLabel{#2}%
    \begin{figure}%
    \centering
}{%
    \caption{\figsCaption}%
    \label{\figsLabel}%
    \end{figure}%
}

% figure envs
\newcommand{\LoadBandFigure}[4][!htp]{
	\begin{figure}[#1]
		\centering
		\includegraphics[width=\columnwidth]{plots/out/load-band-#2.pdf}
		\caption{#3}
		\label{#4}
	\end{figure}
}

% other stuff

\newcommand*\var[1]{\textit{#1}}
\newcommand*\class[1]{\texttt{#1}}
\newcommand*\func[1]{\texttt{#1()}}

%%%%%%%%%%%%%%%%%%%%%%%%%%%%%%%%%%%%%%%%%%%%%%%%%%

\theoremstyle{plain}
\newtheorem{theorem}{Theorem}
\newtheorem{lemma}[theorem]{Lemma}
\newtheorem{corollary}[theorem]{Corollary}

%\newcommand*{\theoremautorefname}{Theorem}
\newcounter{metatestcase}
\setcounter{metatestcase}{0}

\newcounter{testcase}
\setcounter{testcase}{0}

\newcounter{subtestcase}[testcase]
\setcounter{subtestcase}{0}

\newcounter{testcaseresult}
\setcounter{testcaseresult}{0}

\renewcommand*{\thetestcase}{%
    \arabic{metatestcase}.\arabic{testcase}%
}
\renewcommand*{\thesubtestcase}{%
    \thetestcase.\arabic{subtestcase}%
}
\renewcommand*{\theHtestcase}{%
    \arabic{metatestcase}.\arabic{testcase}%
}
\renewcommand*{\theHsubtestcase}{%
    \theHtestcase.\arabic{subtestcase}%
}

\newcommand{\testcaseautorefname}{Testfall}
\newcommand{\subtestcaseautorefname}{Test}

\newcommand{\MetaTestCase}{%
    \stepcounter{metatestcase}%
    \setcounter{testcase}{0}%
    \setcounter{testcaseresult}{0}%
}
\newcommand*{\TestCase}[1]{%
    \ifnum\value{testcase}>0%
        \ifnum\value{testcaseresult}=0%
            \todo{Results TC \thetestcase}%
        \fi%
    \fi%
    \refstepcounter{testcase}%
    \setcounter{subtestcase}{0}%
    \setcounter{testcaseresult}{0}%
    \subsubsection*{Testfall \thetestcase%
        \ifx#1\relax\else: {#1}\fi%
    }%
    \addcontentsline{toc}{subsubsection}{%
        Testfall \thetestcase: {#1}%
    }%
}
\newcommand{\SubTestCase}[1][\relax]{%
    \refstepcounter{subtestcase}%
    \paragraph*{Test \thesubtestcase%
        \ifx#1\relax\else: {#1}\fi%
    }%
}

\newcommand{\FinalizeTestCases}{%
    \ifnum\value{testcase}>0%
        \ifnum\value{testcaseresult}=0%
            \todo{Results TC \thetestcase}%
        \fi%
    \fi%
}

\newcommand{\TestCaseParameter}[2]{\item[#1] {#2}}
\newcommand{\TestCaseMeasuredSize}[2]{\item[#1] {#2}}

\newenvironment{TestCaseDescription}{%
    \paragraph*{Beschreibung}%
}{}
\newenvironment{TestCaseParameters}{%
    \paragraph*{Parameter}%
    \begin{description}%
}{%
    \end{description}%
}
\newenvironment{TestCaseMeasuredSizes}{%
    \paragraph*{Messgrößen}%
    \begin{description}%
}{%
    \end{description}%
}
\newenvironment{TestCaseExpectation}{%
    \paragraph*{Erwartungshaltung}%
}{}
\newenvironment{TestCaseResult}{%
    \stepcounter{testcaseresult}%
    \paragraph*{Auswertung}%
}{}

%%%%%%%%%%%%%%%%%%%%%%%%%%%%%%%%%%%%%%%%%%%%%%%%%%

\colorlet{darkgreen}{green!50!black}
\colorlet{mypurple}{blue!50!red}
\colorlet{b1r5}{blue}
\colorlet{b2r5}{blue!75!red}
\colorlet{b3r5}{blue!50!red}
\colorlet{b4r5}{blue!25!red}
\colorlet{b5r5}{red}

\colorlet{b1r8}{red!0!blue}
\colorlet{b2r8}{red!14!blue}
\colorlet{b3r8}{red!28!blue}
\colorlet{b4r8}{red!42!blue}
\colorlet{b5r8}{red!57!blue}
\colorlet{b6r8}{red!71!blue}
\colorlet{b7r8}{red!85!blue}
\colorlet{b8r8}{red!100!blue}

\pgfplotsset{
	every boxplot/.style={
		mark=o,
		every mark/.append style={
			mark size=0.8pt,
		},
	},
}

% pgfplots stuff
\pgfplotscreateplotcyclelist{custom trans3}{
	blue\\
	blue!50!red\\
	red\\
}

\pgfplotscreateplotcyclelist{custom trans3 fill}{
	blue,fill=blue!25!white\\
	blue!50!red,fill=b3r5!25!white\\
	red,fill=red!25!white\\
}

\pgfplotscreateplotcyclelist{custom trans5}{
	blue\\
	blue!75!red\\
	blue!50!red\\
	blue!25!red\\
	red\\
}

\pgfplotscreateplotcyclelist{custom trans5 fill}{
	blue,fill=b1r5!25!white\\
	blue!75!red,fill=b2r5!25!white\\
	blue!50!red,fill=b3r5!25!white\\
	blue!25!red,fill=b4r5!25!white\\
	red,fill=b5r5!25!white\\
}

\pgfplotscreateplotcyclelist{custom trans4}{
	b1r8\\
	b3r8\\
	b6r8\\
	b8r8\\
}

\pgfplotscreateplotcyclelist{custom trans8}{
	b1r8\\
	b2r8\\
	b3r8\\
	b4r8\\
	b5r8\\
	b6r8\\
	b7r8\\
	b8r8\\
}

\pgfplotscreateplotcyclelist{custom trans8 fill}{
	b1r8,fill=b1r8!25!white\\
	b2r8,fill=b2r8!25!white\\
	b3r8,fill=b3r8!25!white\\
	b4r8,fill=b4r8!25!white\\
	b5r8,fill=b5r8!25!white\\
	b6r8,fill=b6r8!25!white\\
	b7r8,fill=b7r8!25!white\\
	b8r8,fill=b8r8!25!white\\
}

\pgfplotscreateplotcyclelist{custom trans6}{
	b1r8\\
	b2r8\\
	b4r8\\
	b5r8\\
	b7r8\\
	b8r8\\
}

\pgfplotscreateplotcyclelist{custom trans6 fill}{
	b1r8,fill=b1r8!25!white\\
	b2r8,fill=b2r8!25!white\\
	b4r8,fill=b4r8!25!white\\
	b5r8,fill=b5r8!25!white\\
	b7r8,fill=b7r8!25!white\\
	b8r8,fill=b8r8!25!white\\
}

\pgfplotscreateplotcyclelist{custom trans8 box}{
	b1r8,mark=x\\
	b2r8,mark=x\\
	b3r8,mark=x\\
	b4r8,mark=x\\
	b5r8,mark=x\\
	b6r8,mark=x\\
	b7r8,mark=x\\
	b8r8,mark=x\\
}

\newcommand*{\BoxplotElfiThroughputTicks}{1, 2, 3}
\newcommand*{\BoxplotElfiThroughputColumn}{2}
\newcommand{\BoxplotElfiThroughput}[3][0, 10, 30]{
	\begin{tikzpicture}
		\begin{axis}[
			title={RX Throughput in Mbit/s},
			xlabel={Insert Probability in \(\%\)},
			boxplot,
			boxplot/draw direction=y,
			xtick={\BoxplotElfiThroughputTicks},
			xticklabels={#1},
			scaled ticks=false,
		]
			\foreach \i in {#1}{
				\addplot table[y index=\BoxplotElfiThroughputColumn]{stats/out/throughput-stat/elfi.#2.\i.#3-bps.csv};
			}
		\end{axis}
	\end{tikzpicture}
	\begin{tikzpicture}
		\begin{axis}[
			title={RX Throughput in Packets/s},
			xlabel={Insert Probability in \(\%\)},
			boxplot,
			boxplot/draw direction=y,
			xtick={\BoxplotElfiThroughputTicks},
			xticklabels={#1},
			scaled ticks=false,
		]
			\foreach \i in {#1}{
				\addplot table[y index=\BoxplotElfiThroughputColumn]{stats/out/throughput-stat/elfi.#2.\i.#3-pps.csv};
			}
		\end{axis}
	\end{tikzpicture}
}

\newcommand{\PlotImbalance}[1]{
	\begin{tikzpicture}
		\begin{axis}[
			xlabel={Time in s},
			ylabel={Imbalance},
		]
			\addplot[blue] table {stats/out/util-stat/#1-rtb-util-imbalance.csv};
			\addlegendentry{Encryption path}
			\addplot[red]  table {stats/out/util-stat/#1-btr-util-imbalance.csv};
			\addlegendentry{Decryption path}
		\end{axis}
	\end{tikzpicture}
}

\newcommand{\PlotBand}[6][blue]{
	\addplot[#1, name path=L] table[x=#3,y=#5] {#2};
	\addplot[#1, name path=U, forget plot] table[x=#3,y=#6] {#2};
	\addplot[#1!25!white, forget plot] fill between [of=L and U];

	\addplot[black, thick] table[x=#3,y=#4] {#2};
}

\newenvironment{PlotLoadTimeseriesAxis}[1][]{
	\begin{tikzpicture}
		\begin{axis}[
			xlabel={Time [s]},
			ylabel={\(u\)},
			width=11cm,
			height=5cm,
			ymin=-50,
			ymax=1050,
			#1
		]
}{
		\end{axis}
	\end{tikzpicture}
}

\newcommand{\PlotLoadTimeseriesAddplots}[3][]{
	\foreach \i in {#2}{
		\addplot+[
			mark=,
			#1
		] table [
			x=elapsed,
			y=\i
		] {#3};
	}
}

\newenvironment{PlotLoadBandAxis}[1][]{
	\begin{tikzpicture}
		\begin{axis}[
			xlabel={Time [s]},
			ylabel={Utilization},
			width=11cm,
			height=7cm,
			#1
		]
}{
		\end{axis}
	\end{tikzpicture}
}

\newcommand{\PlotLoadBand}[2][blue]{
	\PlotBand[#1]{#2}{elapsed}{average}{average_minus_stddev}{average_plus_stddev}
}
\newcommand{\PlotMinMax}[2][]{
	\addplot+[
		dashed,
		black,
		mark=,
		forget plot,
		#1
	] table [
		x=elapsed,
		y=min,
	] {#2};
	\addplot+[
		dashed,
		black,
		mark=,
		#1
	] table [
		x=elapsed,
		y=max,
	] {#2};
}

\newcommand{\PlotFullLoadBand}[1]{
	\PlotLoadBand{#1}
	\addlegendentry{\(\overline{u} \pm \sigma_u\)}
	\addlegendentry{\(\overline{u}\)}
	\PlotMinMax{#1}
	\addlegendentry{\(u_{min}\)/\(u_{max}\)}
}

\newenvironment{PlotLoadBoxAxis}[1][]{
	\begin{tikzpicture}
		\begin{axis}[
			boxplot/draw direction=y,
			xtick distance=1,
			xlabel={\(C_{RSS}\)},
			ylabel={\(\overline{u}\)},
			extra y tick style={
				grid=major,
			},
			major grid style={
				dashed,
			},
			width=6cm,
			height=6cm,
			ymin=-50,
			ymax=1050,
			#1
		]
}{
		\end{axis}
	\end{tikzpicture}
}

\newcommand{\PlotLoadBox}[2][]{
	\addplot+[
		boxplot,
		#1
	] table [
		x=elapsed,
		y=average,
	] {#2};
}

\newcommand{\PlotLoadMaxBox}[2][]{
	\addplot+[
		boxplot,
		#1
	] table [
		x=elapsed,
		y=max,
	] {#2};
}

\newenvironment{PlotThroughputAxis}[2][]{
	\begin{tikzpicture}
		\begin{axis}[
			boxplot/draw direction=y,
			ylabel={%
				\ifthenelse{\equal{#2}{pps}}{%
					\(t_p\) [\(Pkt/s\)]%
				}{%
					\ifthenelse{\equal{#2}{pps-largest}}{%
						\(t_{{l}_p}\) [\(Pkt/s\)]%
					}{%
						\ifthenelse{\equal{#2}{Mbps-largest}}{%
							\(t_{l_b}\) [\(Mbit/s\)]%
						}{%
							\(t_b\) [\(Mbit/s\)]%
						}%
					}%
				}%
			},
			scaled ticks=false,
			ymin=0,
			width=7cm,
			height=5cm,
			#1
		]
		\ifthenelse{\equal{#2}{pps}}{
			\edef\ba@tpcol{RX-pps}
		}{
			\ifthenelse{\equal{#2}{pps-largest}}{
				\edef\ba@tpcol{RX-pps-first}
			}{
				\ifthenelse{\equal{#2}{Mbps-largest}}{
					\edef\ba@tpcol{RX-Mbps-first}
				}{
					\edef\ba@tpcol{RX-Mbps}
				}
			}
		}
}{
		\end{axis}
	\end{tikzpicture}
}

\newcommand{\PlotThroughputAddplot}[2][]{
	\addplot+ [
		boxplot,
		#1
	] table [
		x=run,
		y=\ba@tpcol,
	] {#2};
}

\newenvironment{PlotLatencyAxis}[2][]{
	\begin{tikzpicture}
		\begin{axis}[
			boxplot/draw direction=y,
			ylabel={\ifthenelse{\equal{#2}{avg}}{\(\overline{l}\)}{\ifthenelse{\equal{#2}{max}}{\(l_{max}\)}{\(l_{min,max}\)}} [\(\mu s\)]},
			scaled ticks=false,
			width=7cm,
			height=5cm,
			ymin=0,
			#1
		]
}{
		\end{axis}
	\end{tikzpicture}
}

\newcommand{\PlotLatencyAddplotAvg}[2][]{
	\addplot+ [
		boxplot,
		#1
	] table [
		x=run,
		y=Delay-avg,
	] {#2};
}

\newcommand{\PlotLatencyAddplotMax}[2][]{
	\addplot+ [
		boxplot,
		#1
	] table [
		x=run,
		y=Delay-max,
	] {#2};
}

\newcommand{\PlotLatencyAddplotMinMax}[3][]{
	\addplot+ [
		boxplot,
		boxplot/draw position=#2,
		#1
	] table [
		x=run,
		y=Delay-min,
	] {#3};
	\addplot+ [
		boxplot,
		boxplot/draw position=#2,
		#1
	] table [
		x=run,
		y=Delay-max,
	] {#3};
}

%\newcommand{\PlotLoadBandYmin}{-50}
%\newcommand{\PlotLoadBandYmax}{1050}
%\newcommand{\PlotLoadBandNode}{l6-rtb}

%\newcommand{\PlotLoadBand}[1]{
%	\begin{tikzpicture}
%		\begin{axis}[
%			xlabel={Time [s]},
%			ylabel={Load},
%			ymin=\PlotLoadBandYmin,
%			ymax=\PlotLoadBandYmax,
%			title={Average \(\pm \sigma\)},
%		]
%			\PlotBand{stats/out/util-stat/#1-\PlotLoadBandNode-util-avg-band.csv}{elapsed}{mean}{lower_band}{upper_band}
%		\end{axis}
%	\end{tikzpicture}
%	\begin{tikzpicture}
%		\begin{axis}[
%			xlabel={Time [s]},
%			ylabel={Load},
%			ymin=\PlotLoadBandYmin,
%			ymax=\PlotLoadBandYmax,
%			title={Median, Quartiles, Minimum and Maximum},
%		]
%			\PlotBand[red]{stats/out/util-stat/#1-\PlotLoadBandNode-util-agg-stat.csv}{elapsed}{median}{perc25_0}{perc75_0}
%			\addplot[dashed, black] table[x=elapsed,y=min] {stats/out/util-stat/#1-\PlotLoadBandNode-util-agg-stat.csv};
%			\addplot[dashed, black] table[x=elapsed,y=max] {stats/out/util-stat/#1-\PlotLoadBandNode-util-agg-stat.csv};
%		\end{axis}
%	\end{tikzpicture}
%}

\newenvironment{PlotIsolatedFlowsAxis}[1][]{
	\begin{tikzpicture}
		\begin{axis}[
			xlabel={Time [s]},
			ylabel={Number of isolated flows},
			xmin=-50,
			height=5cm,
			width=13cm,
			extra x tick style={
				grid=major,
				tick label style={
					anchor=north west,
					font={\small},
					yshift=-1.0em,
				},
			},
			#1
		]
}{
		\end{axis}
	\end{tikzpicture}
}
\newenvironment{PlotIsolatedFlowsAxisRelative}[1][]{
	\begin{tikzpicture}
		\begin{axis}[
			xlabel={Time [s]},
			ylabel={Isolated flows [\%]},
			xmin=-50,
			height=5cm,
			width=13cm,
			extra x tick style={
				grid=major,
				tick label style={
					anchor=north west,
					font={\small},
					yshift=-1.0em,
				},
			},
			legend style={
				cells={anchor=east},
				legend pos=north west,
			},
			#1
		]
}{
		\end{axis}
	\end{tikzpicture}
}
\newenvironment{PlotIsolatedFlowsBoxAxis}[1][]{
	\begin{tikzpicture}
		\begin{axis}[
			xlabel={Skew [\(\%\)]},
			xtick={1,2,3,4},
			xticklabels={0,30,50,90},
			xtick distance=1,
			ylabel={\(I\) [1]},
			bar width=6pt,
			ybar=1pt,
			ymin=0,
			xmin=0.5,
			xmax=4.5,
			height=5cm,
			width=7cm,
			legend style={
				legend pos=north west,
			},
			#1
		]
}{
		\end{axis}
	\end{tikzpicture}
}
\newenvironment{PlotIsolatedFlowsBoxAxisLargest}[1][]{
	\begin{tikzpicture}
		\begin{axis}[
			xlabel={Skew [\(\%\)]},
			xtick={1,2,3,4},
			xticklabels={0,30,50,90},
			xtick distance=1,
			ylabel={\(i_{l}\) [\(\%\)]},
			bar width=6pt,
			xmin=0.5,
			xmax=4.5,
			ybar=1pt,
			ymin=0,
			height=5cm,
			width=7cm,
			legend style={
				legend pos=north west,
			},
			#1
		]
}{
		\end{axis}
	\end{tikzpicture}
}
\newcommand{\PlotIsolatedFlowsAddplot}[2][]{
	\addplot+[
		sharp plot,
		mark=,
		#1
	]
	table [
		x=elapsed,
		y=elfi-isolated,
	] {#2};
}
\newcommand{\PlotIsolatedFlowsAddplotRelative}[2][]{
	\addplot+[
		sharp plot,
		mark=,
		#1
	]
	table [
		x=elapsed,
		y=elfi-isolated-percent,
	] {#2};
}
\newcommand{\PlotIsolatedFlowsAddplotBox}[2][]{
	\addplot+[
		ybar,
		mark=,
		#1
	] table [
		x=run,
		y=elfi-isolated-avg,
	]{#2};
}
\newcommand{\PlotIsolatedFlowsAddplotBoxLargest}[2][]{
	\addplot+[
		ybar,
		mark=,
		#1
	] table [
		x=run,
		y=elfi-isolated-first-percent,
	]{#2};
}

\newenvironment{PlotIsolatedThroughputAxis}[1][\relax]{
	\begin{tikzpicture}
		\begin{axis}[
			xtick={1,2,3,4},
			xticklabels={0,30,50,90},
			xlabel={Skew [\(\%\)]},
			ylabel={Isolated Throughput [\(Pkt/s\)]},
			bar width=6pt,
			ybar=1pt,
			width=7cm,
			height=5cm,
			ymin=0,
			#1
		]
}{
		\end{axis}
	\end{tikzpicture}
}

\newcommand{\PlotIsolatedThroughputAddplot}[2][]{
	\addplot+[
		ybar,
		#1
	] table [
		x={run},
		y={elfi-isolated-pps},
	] {#2};
}

\newcommand{\PlotIsolatedFlows}[1]{
	\begin{PlotIsolatedFlowsAxis}
		\PlotIsolatedFlowsAddplot{#1}
	\end{PlotIsolatedFlowsAxis}
}

\makeatother

\usepackage[acronym]{glossaries}
\makeglossaries
\renewcommand{\glossarysection}[2][]{}
\usepackage{pifont}
\newcommand{\uncheckmark}{\ding{55}}
\newcommand{\unchecked}{---}

\newcommand{\req}[2] {%
\parbox[t]{\linewidth}{% 
\makebox[3.5em][l]{#1}%
 \parbox[t]{\dimexpr\linewidth-3.5em\relax}{#2}% 
 }% 
 }

 \usepackage{rotating}